% SPDX-License-Identifier: PMPL-1.0-or-later
% SPDX-FileCopyrightText: 2026 Jonathan D.A. Jewell <j.d.a.jewell@open.ac.uk>
%
% arXiv-style academic paper on emotional lineage licensing
% and consent layers for non-interpretive AI systems.

\documentclass[11pt,twocolumn]{article}

% --- Packages ---
\usepackage[utf8]{inputenc}
\usepackage[T1]{fontenc}
\usepackage{lmodern}
\usepackage[margin=2cm]{geometry}
\usepackage{hyperref}
\usepackage{cite}
\usepackage{amsmath,amssymb}
\usepackage{graphicx}
\usepackage{booktabs}
\usepackage{enumitem}
\usepackage{xcolor}
\usepackage{listings}
\usepackage{fancyhdr}
\usepackage{titlesec}
\usepackage{abstract}
\usepackage{microtype}
\usepackage{url}

% --- Hyperlink styling ---
\hypersetup{
    colorlinks=true,
    linkcolor=blue!60!black,
    citecolor=blue!60!black,
    urlcolor=blue!60!black
}

% --- Listing style ---
\lstset{
    basicstyle=\ttfamily\footnotesize,
    breaklines=true,
    frame=single,
    backgroundcolor=\color{gray!10},
    xleftmargin=1em,
    xrightmargin=1em
}

% --- Title ---
\title{%
    \textbf{Emotional Lineage and Consent Layers:\\
    Extending Open-Source Licensing for the\\
    Age of Non-Interpretive AI Systems}
}

\author{%
    Jonathan D.A.\ Jewell\\
    \textit{Independent Researcher}\\
    \texttt{j.d.a.jewell@open.ac.uk}\\
    \texttt{https://github.com/hyperpolymath}
}

\date{March 2026}

\begin{document}

\maketitle

% ============================================================================
% ABSTRACT
% ============================================================================
\begin{abstract}
Open-source software licenses were designed in an era when the primary
consumers of source code were human developers. The four freedoms of free
software---to use, study, share, and improve---implicitly assume a human
agent capable of \emph{interpreting} the work: understanding its purpose,
respecting its context, and engaging with its meaning. The rise of
large-scale non-interpretive AI systems---those that ingest, compress, and
recombine creative works without semantic comprehension---exposes a
structural gap in every major open-source license. Neither the MIT License,
the GNU General Public License, the Apache License 2.0, nor the Creative
Commons family addresses the distinction between a human reading code and
a statistical model consuming it as training data. This paper introduces
two complementary concepts developed for the Palimpsest-MPL License
(PMPL-1.0): \emph{emotional lineage protection}, which formalises the
obligation to preserve the narrative, cultural, and symbolic meaning
embedded in creative works; and the \emph{consent layer for
non-interpretive systems}, which establishes an explicit opt-in mechanism
for automated ingestion of licensed works. We present formal definitions,
examine compatibility with existing legal frameworks including the EU AI
Act and U.S.\ fair use doctrine, propose a machine-readable consent
metadata standard, and argue that the interpretive/non-interpretive
distinction represents a necessary evolution in software licensing theory.
The paper concludes that consent-based licensing, far from being anti-AI,
is the natural extension of the principles that built the open-source
movement.
\end{abstract}

\medskip
\noindent\textbf{Keywords:} software licensing, open source, AI training,
emotional lineage, consent, non-interpretive systems, copyleft, PMPL,
cultural preservation, SPDX

% ============================================================================
% 1. INTRODUCTION
% ============================================================================
\section{Introduction}
\label{sec:introduction}

The open-source software movement transformed how humanity builds digital
infrastructure. From the GNU Manifesto~\cite{stallman1985gnu} to the
proliferation of permissive licenses in the 2010s, the legal instruments
governing source code have co-evolved with the communities that produce
it. Yet these instruments share a common, unstated assumption: that the
entity exercising the granted rights is a \emph{human being}, or at
minimum, an organisation directed by human judgment.

This assumption was invisible because it was universal. When Richard
Stallman articulated the four freedoms~\cite{stallman2002free}, the idea
that a non-human system might ``use'' software---not as a tool, but as
raw material for pattern extraction---was science fiction. When the MIT
License declared that ``permission is hereby granted\ldots to any person
obtaining a copy,'' the word ``person'' was not legally contested. When
the Apache Software Foundation drafted patent clauses, the threat model
was corporate litigation, not statistical regurgitation.

The emergence of large language models (LLMs), diffusion models, and
other generative AI systems has rendered these assumptions untenable.
These systems do not \emph{use} software in any sense contemplated by
existing licenses. They do not read documentation, understand variable
names, or appreciate the political context of a commit message. They
consume bytes. They extract statistical regularities. They produce
outputs that may superficially resemble---but do not \emph{understand}---%
the works they were trained on.

This distinction between \emph{interpretive use} (a human engaging with
a work's meaning) and \emph{non-interpretive use} (an automated system
processing a work as data) is the central concern of this paper. We argue
that it represents not merely a technical nuance, but a \emph{category
error} in the foundations of open-source licensing---one that, left
unaddressed, threatens both creator autonomy and the sustainability of
open-source ecosystems.

The paper proceeds as follows. Section~\ref{sec:background} surveys the
history of open-source licensing and the current legal landscape around
AI training. Section~\ref{sec:emotional-lineage} introduces the concept
of emotional lineage and argues for its formal recognition in licensing
frameworks. Section~\ref{sec:interpretive-distinction} develops the
interpretive/non-interpretive distinction with formal definitions and
legal grounding. Section~\ref{sec:consent-layer} presents the consent
layer mechanism. Section~\ref{sec:comparison} compares the approach with
existing licenses. Section~\ref{sec:palimpsest} explains the palimpsest
metaphor and its implications. Section~\ref{sec:legal-analysis} analyses
enforceability. Section~\ref{sec:implementation} describes the
machine-readable implementation. Section~\ref{sec:discussion} discusses
broader implications, and Section~\ref{sec:conclusion} concludes.

% ============================================================================
% 2. BACKGROUND
% ============================================================================
\section{Background}
\label{sec:background}

\subsection{A Brief History of Open-Source Licensing}

The legal infrastructure of open source emerged in three waves. The
\emph{first wave} (1983--1998) was dominated by the GNU General Public
License (GPL), which established the copyleft principle: derivative works
must be distributed under the same terms~\cite{stallman2002free}. The
GPL's genius was recognising that freedom, without legal mechanism, is
merely aspiration. Its limitation was treating software as pure
\emph{functionality}---something compiled and executed, not something
carrying cultural weight.

The \emph{second wave} (1998--2015) saw the rise of permissive licenses.
The MIT License, the BSD licenses, and the Apache License 2.0 reflected
a pragmatic turn: maximum adoption through minimum obligation. The Open
Source Initiative (OSI) codified the Open Source Definition, which
deliberately excluded ethical considerations from licensing
criteria~\cite{perens1999osi}. This was a philosophical choice, not an
oversight. The OSI position was that licenses should govern
\emph{distribution mechanics}, not \emph{use}.

The \emph{third wave} (2015--present) has been characterised by tension.
The rise of cloud computing produced the Server Side Public License
(SSPL) and the Business Source License (BSL), which attempted to close
the ``SaaS loophole''~\cite{mongodb2018sspl}. The Hippocratic
License~\cite{ehmke2019hippocratic} and the Anti-996
License~\cite{anti996} introduced ethical constraints, provoking fierce
debate about whether such constraints were compatible with ``open source''
as defined by the OSI. Simultaneously, the AI training problem emerged
as a distinct challenge, largely unaddressed by any existing licence.

\subsection{The AI Training Problem}

When a large language model is trained on a corpus that includes
GPL-licensed code, does the resulting model constitute a ``derivative
work''? The question is not merely academic. GitHub Copilot, trained
partly on GPL-licensed repositories, can reproduce substantial portions
of copylefted code~\cite{copilot2022controversy}. If the model is a
derivative work, it arguably must be distributed under GPL terms---an
obligation that no model provider has accepted. If it is not, then
copyleft has a gap wide enough to drive an entire industry through.

The legal landscape is unsettled. In \emph{Andersen v.\ Stability AI}
(2023), artists challenged the use of their work in training diffusion
models~\cite{andersen2023}. The \emph{New York Times v.\ OpenAI} (2023)
lawsuit raised similar questions for text~\cite{nyt2023openai}. Courts
have not yet produced definitive rulings, and the answers may differ by
jurisdiction.

What is clear is that existing licenses were not designed to answer these
questions. The GPL speaks of ``source code'' and ``object code,'' of
``conveying'' and ``distribution''~\cite{gplv3}. It does not speak of
training data, model weights, or statistical reproduction. The MIT
License grants permission to ``use, copy, modify, merge, publish,
distribute, sublicense, and/or sell copies''---a list that does not
obviously include ``feed to a neural network as training
data''~\cite{mit-license}. The Apache License 2.0 includes patent
grants but says nothing about whether ingestion into a model constitutes
``use''~\cite{apache2}.

\subsection{The EU AI Act and Emerging Regulation}

The European Union's AI Act~\cite{eu-ai-act-2024}, which entered into
force in 2024, represents the most comprehensive regulatory framework
for AI systems to date. Article 53 requires providers of general-purpose
AI models to ``put in place a policy to comply with Union copyright law,''
including ``the identification and compliance with reservations of rights
expressed pursuant to Article 4(3) of Directive (EU) 2019/790''---the
text and data mining opt-out provision.

This is significant because it establishes, at the regulatory level, the
principle that \emph{rightsholders may opt out of AI training}. The EU
AI Act does not create the right---it presupposes it. The question for
licensing is how to \emph{exercise} that right in a machine-readable,
enforceable manner.

\subsection{The Limits of Copyright}

Copyright law, as traditionally understood, protects the \emph{expression}
of ideas, not the ideas themselves~\cite{samuelson1984copyright}.
This distinction, known as the idea-expression dichotomy, creates a
structural problem for creative works whose value lies not in their
literal expression but in their \emph{meaning}. A protest song is
valuable not because of its specific chord progressions (expression) but
because of what it represents---the movement it accompanied, the
injustice it named, the community it galvanised.

Copyright can prevent someone from copying the chord progressions. It
cannot prevent someone from stripping the song of its context, feeding
it to a generative model, and producing ``content'' that
statistically resembles protest music without carrying any of its
purpose. This is the gap that emotional lineage protection is designed
to address.

% ============================================================================
% 3. THE EMOTIONAL LINEAGE CONCEPT
% ============================================================================
\section{The Emotional Lineage Concept}
\label{sec:emotional-lineage}

\subsection{Definition}

We define \textbf{emotional lineage} as:

\begin{quote}
\textit{The narrative, cultural, symbolic, and contextual meaning embedded
in creative works, including but not limited to: protest traditions,
cultural heritage, trauma narratives, community stories, and the social
relationships that produced the work.}
\end{quote}

This definition is deliberately broad. It encompasses the historical
context of a codebase (why it was written, what problem it emerged from),
the cultural milieu of a creative work (the tradition it belongs to, the
community it serves), and the relational fabric of collaborative
production (the human connections, mentorships, and conflicts that shaped
its development).

\subsection{Why Code Has Emotional Lineage}

The claim that source code carries emotional lineage may seem
counterintuitive to those who view programming as a purely technical
activity. We argue this view is mistaken, on several grounds.

\textbf{First, code embeds values.} The choice to build an encryption
library reflects a belief in privacy. The decision to make a medical
records system open-source reflects a commitment to healthcare access.
The architecture of a decentralised protocol reflects a political
philosophy. These are not incidental; they are constitutive of the
software's identity.

\textbf{Second, code carries community history.} A project's git log is
a social document. It records who contributed, when, and why. It shows
debates in commit messages, struggles in issue trackers, celebrations in
release notes. When this history is compressed into model weights, the
social fabric is destroyed while the statistical patterns are preserved.

\textbf{Third, code exists within traditions.} The Unix philosophy, the
Lisp tradition, the Haskell community's emphasis on type safety, the
Rust community's focus on safety without garbage collection---these are
not merely technical preferences. They are \emph{cultures}, with values,
norms, heroes, and histories. Code written within these traditions
carries their meaning.

\textbf{Fourth, some code is explicitly political.} Encryption tools
built for activists, anonymity networks designed for whistleblowers,
accessibility libraries created for disabled users---these works were
produced not for abstract ``utility'' but for specific human purposes,
often at personal risk to their creators. To strip this context is not
a neutral act.

\subsection{Cultural Erasure Through Context Collapse}

When a non-interpretive system ingests a creative work, it performs what
we term \textbf{context collapse}: the reduction of a multi-layered,
historically situated work to a set of statistical features. This is not
merely a loss of attribution (which existing licenses already address)
but a loss of \emph{meaning}.

Consider an analogy from cultural heritage. A sacred object in a museum
retains its physical form but loses its sacred function. The museum
visitor sees an artefact; the originating community sees a desecration.
The object's ``emotional lineage''---its relationship to the community
that created it, the ceremonies it was part of, the beliefs it
embodies---has been severed~\cite{clifford1988predicament}.

The same dynamic occurs when AI systems ingest creative works. The
\emph{expression} is statistically preserved (the model can reproduce
patterns that resemble the original). The \emph{meaning} is destroyed
(the model has no access to purpose, context, or significance). This
is not an accident of current technology that better AI will solve. It
is a structural feature of non-interpretive processing.

\subsection{Formalising Preservation Obligations}

The Palimpsest-MPL License (PMPL-1.0) addresses emotional lineage
through Section 3.1, which requires that distributors and derivative
work creators ``make reasonable efforts to preserve and communicate the
Emotional Lineage of Covered Software\ldots including maintaining
narrative context, cultural attributions, and symbolic meaning where
documented.''

The ``where documented'' qualifier is important. The obligation is not
to divine unspoken intentions, but to preserve context that has been
\emph{explicitly recorded}. This is operationalised through several
mechanisms:

\begin{enumerate}[label=(\alph*)]
    \item A \texttt{LINEAGE.md} file that documents the origin story,
          cultural context, and symbolic meaning of the work.
    \item Provenance metadata that cryptographically links contributions
          to their authors and timestamps.
    \item A prohibition on stripping or obscuring provenance metadata
          from distributed works (Section 4.1).
\end{enumerate}

The standard of ``reasonable efforts'' is borrowed from existing legal
frameworks (e.g., the ``reasonable person'' standard in tort law) and
deliberately avoids imposing disproportionate burdens. A developer
forking a project is not required to write an essay on its cultural
significance. They \emph{are} required not to delete the essay if one
exists.

% ============================================================================
% 4. THE INTERPRETIVE / NON-INTERPRETIVE DISTINCTION
% ============================================================================
\section{The Interpretive/Non-Interpretive Distinction}
\label{sec:interpretive-distinction}

\subsection{Formal Definition}

We distinguish two categories of use:

\begin{description}[style=nextline]
    \item[\textbf{Interpretive use}]
    Use by an agent (human or system) that engages with the semantic
    content of a work---understanding its purpose, respecting its
    context, and producing outputs that reflect comprehension of its
    meaning. Examples include: a developer reading and modifying source
    code; a musician listening to and building upon a composition; a
    student studying a textbook.

    \item[\textbf{Non-interpretive use}]
    Use by an automated system that processes a work without semantic
    engagement---extracting statistical patterns, features, or structures
    without comprehending purpose, context, or meaning. Examples include:
    AI model training on source code; content aggregation systems that
    scrape and recombine text; automated summarisation tools that compress
    meaning into statistical approximations.
\end{description}

The distinction is not about the \emph{sophistication} of the system. A
simple grep command is interpretive in our framework if directed by a
human for a comprehension purpose. A highly sophisticated LLM is
non-interpretive if it processes works as training data without engaging
with their meaning. The criterion is \emph{semantic engagement}, not
computational complexity.

\subsection{The Photography Analogy}

This distinction has precedent in existing law. In privacy and data
protection law, a fundamental difference exists between \emph{viewing}
and \emph{processing}. A person may look at another person's face in a
public space (interpretive engagement). They may not, without consent,
photograph that face and use it to train a facial recognition system
(non-interpretive processing)~\cite{gdpr2016}.

The analogy is precise. Looking at a face involves understanding---you
recognise the person, read their expression, engage with their humanity.
Training a facial recognition system involves extraction---the face
becomes a vector of measurements, stripped of identity, emotion, and
context. The same physical input (photons reflected from a face) is
subject to fundamentally different legal regimes depending on how it is
\emph{used}.

We argue that the same principle should apply to creative works. Reading
code, studying it, learning from it, modifying it---these are
interpretive acts, freely permitted under open-source principles.
Feeding code to a training pipeline that extracts patterns without
comprehension---this is a non-interpretive act, requiring explicit
consent.

\subsection{Why the Distinction Matters}

One might object that the distinction is artificial. After all, human
learning also involves ``pattern extraction'' at a neurological level.
When a developer reads idiomatic Rust code and later writes similar
patterns, have they not ``trained'' on the original?

The objection fails for three reasons.

\textbf{First, human learning is inherently interpretive.} When a
developer reads Rust code, they understand what the borrow checker
enforces and \emph{why}. They grasp the design philosophy. They can
articulate the trade-offs. They form opinions about the code's quality.
This is categorically different from a model that represents
\texttt{fn main()} as a sequence of tokens with statistical
co-occurrence patterns.

\textbf{Second, human learning is bounded by human capacity.} A
developer can read thousands of projects over a career. An LLM can
ingest millions of repositories in hours. The \emph{scale} of
non-interpretive processing creates effects---market displacement,
cultural homogenisation, community disinvestment---that have no analogue
in human learning.

\textbf{Third, human learning produces human agency.} A developer who
learns from open-source code becomes a more capable contributor to the
ecosystem. An AI model trained on open-source code may displace the
developers who wrote it, without contributing to the community that
produced the training data~\cite{lemley2023generative}. The
relationship between consumption and contribution is fundamentally
different.

\subsection{Edge Cases and the Spectrum of Interpretation}

We acknowledge that the boundary between interpretive and
non-interpretive use is not always sharp. Consider the following cases:

\begin{itemize}
    \item \textbf{IDE autocompletion}: A system that suggests code
          completions based on local context. This is borderline:
          the system does not ``understand'' the code, but it is
          directed by a human who does, and operates at human scale.
    \item \textbf{Static analysis}: A tool that detects bugs by
          pattern matching. This is non-interpretive in mechanism
          but interpretive in purpose---it serves human comprehension.
    \item \textbf{Search indexing}: A system that indexes code to
          make it findable. This processes code without understanding
          it, but serves the interpretive use of human searchers.
\end{itemize}

These edge cases suggest that the distinction should be understood as a
\emph{spectrum} rather than a binary. The PMPL-1.0 addresses this by
focusing on the clearest cases---AI training pipelines, content
aggregation systems, and automated summarisation tools---while leaving
the boundary zone to evolve through community practice and, where
necessary, adjudication.

The key principle is: \emph{when in doubt, ask}. The consent layer
(Section~\ref{sec:consent-layer}) provides the mechanism for asking.

% ============================================================================
% 5. THE CONSENT LAYER
% ============================================================================
\section{The Consent Layer}
\label{sec:consent-layer}

\subsection{Design Principles}

The consent layer is governed by four principles:

\begin{enumerate}
    \item \textbf{Default-open for interpretive use.} The license imposes
          no additional restrictions on human use of covered works. The
          four freedoms are fully preserved.
    \item \textbf{Explicit consent for non-interpretive use.} Automated
          systems that process covered works as training data or for
          content aggregation must obtain or verify consent.
    \item \textbf{Granular control.} Creators can specify consent at
          the project level, file level, or even function level.
    \item \textbf{Machine readability.} Consent metadata must be
          processable by automated systems, enabling compliance at scale.
\end{enumerate}

\subsection{Mechanism}

The consent layer operates through three components:

\subsubsection{The Non-Interpretive System Notice}

Section 3.2 of the PMPL-1.0 requires that any entity using covered
software as input to a non-interpretive system must: (a) document such
use in a publicly accessible manner, and (b) not claim that outputs of
such systems carry the emotional lineage of the original work without
explicit permission from contributors.

This is a \emph{transparency obligation}, not a prohibition. The
license does not forbid AI training on covered works. It requires that
such training be \emph{documented}---that the creator and the public
can know when and how their work was used.

\subsubsection{The Ethical Use Declaration}

Section 3.3 requires commercial users to acknowledge that they have read
and understood the Ethical Use Guidelines (Exhibit A) and agree to act in
good faith accordance with their principles. This is modelled on existing
mechanisms in commercial licensing (e.g., Oracle's Java SE License, which
requires acceptance of terms before use).

\subsubsection{Consent Metadata}

The machine-readable component (detailed in Section~\ref{sec:implementation})
allows creators to express their consent preferences in a structured
format. A \texttt{.palimpsest/consent.json} file (or equivalent SPDX
annotation) can specify:

\begin{lstlisting}
{
  "consent": {
    "interpretive": "granted",
    "non-interpretive": {
      "ai-training": "requires-notice",
      "content-aggregation": "denied",
      "search-indexing": "granted",
      "academic-research": "granted"
    }
  },
  "contact": "maintainer@example.org",
  "updated": "2026-01-15"
}
\end{lstlisting}

\subsection{Opt-In vs.\ Opt-Out}

A critical design decision is whether non-interpretive use requires
\emph{opt-in} (explicit grant) or merely respects \emph{opt-out}
(explicit denial). The EU's text and data mining framework under
Directive 2019/790 adopts an opt-out model for commercial purposes: TDM
is permitted unless the rightsholder has ``expressly reserved'' their
rights~\cite{eu-dsm-directive}.

The PMPL-1.0 takes a middle path. Interpretive use is opt-in by default
(the license grants it). Non-interpretive use requires \emph{notice}
(documentation of use) but not necessarily \emph{permission}. Creators
who wish to deny non-interpretive use entirely may do so through the
consent metadata. Creators who wish to permit it may do so, with or
without conditions.

This design reflects a recognition that the open-source ecosystem
encompasses diverse views. Some creators welcome AI training on their
work and see it as a form of dissemination. Others view it as
exploitative extraction. The consent layer does not impose a single
answer; it provides the \emph{mechanism} for each creator to express
their own.

\subsection{Revocability and Versioning}

Consent is versioned and revocable. A creator who grants consent for
AI training in version 1.0 of their work may revoke it in version 2.0.
This does not retroactively invalidate training performed on version 1.0
(which was lawfully obtained under the then-current terms), but it does
prevent future training on version 2.0 without re-obtaining consent.

This is analogous to how GPL version selection works: you may use a work
under the GPL version it was released under, even if the project later
adopts a different license~\cite{gplv3}. Consent, like licensing, is
time-bound to the version.

% ============================================================================
% 6. COMPARISON WITH EXISTING LICENSES
% ============================================================================
\section{Comparison with Existing Licenses}
\label{sec:comparison}

Table~\ref{tab:comparison} summarises how existing licenses address the
concerns raised in this paper.

\begin{table*}[t]
\centering
\caption{Comparison of license provisions relevant to AI and cultural context.}
\label{tab:comparison}
\small
\begin{tabular}{@{}lcccccc@{}}
\toprule
\textbf{Feature} & \textbf{MIT} & \textbf{GPL-3.0} & \textbf{Apache-2.0} & \textbf{MPL-2.0} & \textbf{CC-BY-4.0} & \textbf{PMPL-1.0} \\
\midrule
Attribution required         & \checkmark & \checkmark & \checkmark & \checkmark & \checkmark & \checkmark \\
Copyleft                     &            & \checkmark &            & File-level &            & File-level \\
Patent grant                 &            & \checkmark & \checkmark & \checkmark &            & \checkmark \\
AI training addressed        &            &            &            &            &            & \checkmark \\
Emotional lineage            &            &            &            &            &            & \checkmark \\
Consent mechanism            &            &            &            &            &            & \checkmark \\
Interpretive/non-int.\ dist. &            &            &            &            &            & \checkmark \\
Provenance metadata          &            &            &            &            &            & \checkmark \\
Machine-readable consent     &            &            &            &            &            & \checkmark \\
\bottomrule
\end{tabular}
\end{table*}

\subsection{MIT License}

The MIT License~\cite{mit-license} is maximally permissive. It permits
``any person obtaining a copy'' to ``use, copy, modify, merge, publish,
distribute, sublicense, and/or sell'' the software. The only obligation
is to include the copyright notice and permission notice. AI training
is not mentioned, and the breadth of ``use'' arguably encompasses it.
There is no mechanism for preserving context, lineage, or meaning. The
MIT License treats software as a commodity; the PMPL treats it as a
cultural artefact.

\subsection{GNU General Public License v3}

The GPL-3.0~\cite{gplv3} is the strongest copyleft license in wide use.
Its ``derivative work'' trigger is powerful: if you modify GPL code and
distribute the result, you must release your modifications under the GPL.
However, whether an AI model constitutes a ``derivative work'' of its
training data is deeply contested~\cite{lemley2023generative}. If model
weights are not ``source code'' and inference is not ``distribution,''
the GPL's copyleft trigger never fires. The GPL also says nothing about
emotional lineage, context preservation, or consent for non-interpretive
use.

\subsection{Apache License 2.0}

The Apache License~\cite{apache2} adds explicit patent grants to
permissive licensing. Its patent termination clause (Section 3) provides
protection against patent trolling. However, like the MIT License, it
is silent on AI training, emotional lineage, and non-interpretive use.
Its ``NOTICE'' file requirement is the closest analogue to provenance
preservation, but it does not extend to cultural or narrative context.

\subsection{Mozilla Public License 2.0}

The MPL-2.0~\cite{mpl2} is the direct ancestor of the PMPL-1.0. Its
file-level copyleft strikes a balance between the permissiveness of MIT
and the project-level copyleft of the GPL. However, like all licenses
in this family, it was designed for a world where software is compiled
and distributed by humans, not ingested by statistical models. The PMPL
extends the MPL with the provisions described in this paper while
preserving full MPL-2.0 compatibility.

\subsection{Creative Commons Licenses}

The Creative Commons family, particularly CC-BY-4.0~\cite{cc-by}, is
widely used for non-software creative works. CC-BY requires attribution,
which might seem to address context preservation. However, attribution
in the CC sense is a \emph{credit} obligation, not a \emph{meaning}
obligation. You must say who created the work; you need not say \emph{why}
they created it, what it \emph{means}, or what community it belongs to.
Furthermore, CC licenses do not distinguish between interpretive and
non-interpretive use. A CC-BY photograph can be attributed in an AI
training dataset manifest; this satisfies the license while doing nothing
to preserve the photograph's cultural significance.

\subsection{Data-Specific Licenses}

More recent efforts, such as the Responsible AI License
(RAIL)~\cite{rail2022} and the Open Data Commons licenses, have begun
to address AI-specific concerns. The RAIL framework introduces use
restrictions that are triggered by the \emph{purpose} of use, not just
the \emph{act} of distribution. This is closer to the PMPL approach,
but the RAIL licenses focus on \emph{prohibiting harmful outputs} rather
than \emph{preserving input meaning}. They are complementary rather than
competing frameworks.

% ============================================================================
% 7. THE PALIMPSEST METAPHOR
% ============================================================================
\section{The Palimpsest Metaphor}
\label{sec:palimpsest}

\subsection{Layers of Meaning}

A palimpsest is a manuscript page that has been written on, scraped
clean, and written on again---sometimes multiple times. Medieval scribes,
working with expensive parchment, would erase earlier texts to reuse the
material. But the erasure was never complete. Under ultraviolet light or
multispectral imaging, the earlier layers become visible again. Some of
the most important texts in Western history---Archimedes' \emph{Method
of Mechanical Theorems}, Cicero's \emph{De re publica}---survived only
as palimpsest underlayers~\cite{netz2007archimedes}.

The metaphor is precise. Creative works, including software, are
palimpsests. The visible layer---the current version of the code, the
latest release---sits atop earlier layers: previous versions, abandoned
approaches, design decisions that were debated and resolved, cultural
contexts that shaped choices. These layers are not visible to a
non-interpretive system, which sees only the top layer (the current
bytes). But they are \emph{constitutive} of the work's identity.

\subsection{Preservation Obligations}

The palimpsest metaphor implies a specific ethical obligation: to
preserve \emph{all layers}, not just the top one. This is precisely what
provenance metadata, emotional lineage documentation, and the
\texttt{LINEAGE.md} file are designed to do. They make the underlayers
\emph{accessible}---not by requiring everyone to read them, but by
ensuring they are not \emph{destroyed}.

The distinction matters. A developer forking a project need not study its
complete history. But they must not delete the history. A company using
an open-source library need not understand the cultural context of every
dependency. But they must not strip the metadata that records it. The
obligation is to \emph{not destroy}, not to \emph{fully comprehend}.

This is a lower bar than it might appear. In practice, it means: do not
strip SPDX headers. Do not delete \texttt{LINEAGE.md}. Do not remove
provenance signatures. Do not claim that AI-generated outputs carry the
authority of human-created originals. These are obligations of
\emph{restraint}, not of \emph{effort}.

\subsection{The Palimpsest as Counter-Narrative to ``Move Fast and Break Things''}

The palimpsest metaphor also serves a rhetorical function. It stands in
deliberate opposition to the Silicon Valley ethos of ``move fast and
break things''---an ethos that treats existing structures as obstacles
rather than as repositories of accumulated wisdom. The palimpsest says:
what came before has value. The layers beneath the surface contain
meaning. Speed and disruption are not always virtues; sometimes, the
most important act is preservation.

This is not a Luddite position. The PMPL explicitly permits AI training,
automated processing, and all forms of technological innovation. It
merely insists that innovation should not require the destruction of
context. You can write new text on the parchment; you must not burn the
parchment to fuel the ink.

% ============================================================================
% 8. LEGAL ANALYSIS
% ============================================================================
\section{Legal Analysis}
\label{sec:legal-analysis}

\subsection{Enforceability of Emotional Lineage Provisions}

The enforceability of emotional lineage provisions depends on jurisdiction
and the legal theory under which they are advanced. We consider three
frameworks.

\subsubsection{Contract Law}

Software licenses function as contracts (or, in some jurisdictions, as
bare licenses). The PMPL's emotional lineage provisions are contractual
obligations: by using covered software, you agree to preserve documented
emotional lineage. This is enforceable to the same extent as any other
license term. Courts routinely enforce license conditions that go beyond
copyright's minimum requirements---the GPL's copyleft obligations, for
instance, are contractual additions to copyright's baseline
permissions~\cite{jacobsen2008}.

The ``reasonable efforts'' standard in Section 3.1 is deliberately
chosen for enforceability. Courts are familiar with reasonableness
standards and have extensive precedent for applying
them~\cite{farnsworth2004contracts}. A court need not determine the
``correct'' level of emotional lineage preservation; it need only
determine whether the defendant made \emph{reasonable efforts}.

\subsubsection{Moral Rights}

Many jurisdictions recognise ``moral rights'' (droit moral) that protect
the integrity of creative works independent of economic
rights~\cite{berne1886}. The Berne Convention, to which most nations are
signatories, guarantees the right of attribution and the right to object
to derogatory treatment of a work. Emotional lineage provisions can be
understood as an elaboration of the integrity right: stripping cultural
context from a work is a form of derogatory treatment.

The strength of this argument varies by jurisdiction. France and Germany
have strong moral rights traditions; the United States has weak ones
(limited primarily to visual arts under the Visual Artists Rights Act of
1990)~\cite{vara1990}. The PMPL's contractual approach ensures
enforceability regardless of the local moral rights regime.

\subsubsection{Cultural Heritage Law}

For works with explicit cultural heritage significance, additional legal
frameworks may apply. The UNESCO Convention on Intangible Cultural
Heritage~\cite{unesco2003} and the WIPO Intergovernmental Committee on
Intellectual Property and Genetic Resources, Traditional Knowledge and
Folklore~\cite{wipo-tk} provide frameworks for protecting cultural
expressions. While these instruments do not directly apply to software
licenses, they establish the principle that cultural context is legally
cognisable---a principle that supports the PMPL's approach.

\subsection{Enforceability of the Consent Layer}

The consent layer raises distinct enforceability questions.

\subsubsection{The Notice Obligation}

Section 3.2(a) requires documentation of non-interpretive use. This is
a straightforward contractual obligation analogous to the GPL's source
code disclosure requirement. Failure to document AI training use is a
license violation, with the same remedies available as for any other
license breach (injunction, damages, termination of the license
grant)~\cite{jacobsen2008}.

\subsubsection{The Authority Claim Prohibition}

Section 3.2(b) prohibits claiming that AI outputs carry the emotional
lineage of the original work. This is enforceable under both contract
law (as a license condition) and consumer protection law (as a
prohibition on misleading claims). If a company represents that its
AI-generated code carries the imprimatur of a well-known open-source
project, and that representation is based on training data rather than
actual endorsement, this is straightforwardly actionable as a false or
misleading claim under consumer protection statutes in most
jurisdictions~\cite{ftc-act}.

\subsubsection{Compatibility with Fair Use / Fair Dealing}

In the United States, fair use (17 U.S.C.\ \S 107) provides a defence
to copyright infringement based on four factors: purpose and character
of use, nature of the copyrighted work, amount used, and market
effect~\cite{us-copyright-act}. AI training proponents argue that
training is transformative fair use. Even if courts accept this argument
for copyright purposes, it does not defeat contractual obligations.
Fair use is a defence to \emph{copyright infringement}; it is not a
defence to \emph{breach of contract}. If a user has accepted the PMPL's
terms, the consent layer's obligations apply regardless of whether the
underlying use would constitute fair use in the absence of a
license~\cite{davidson2004click}.

Similarly, in jurisdictions with fair dealing exceptions (UK, Canada,
Australia), the contractual nature of the PMPL's provisions operates
independently of the copyright exception.

\subsection{Interaction with the EU AI Act}

The PMPL's consent layer is designed for compatibility with Article 53 of
the EU AI Act. The machine-readable consent metadata
(Section~\ref{sec:implementation}) provides the ``reservation of rights''
mechanism contemplated by the Act's reference to Directive 2019/790
Article 4(3). AI model providers who comply with the PMPL's consent
layer will, by construction, satisfy the AI Act's transparency and
opt-out compliance requirements for covered works.

\subsection{OSI Compatibility}

The PMPL-1.0 is not OSI-approved, and we anticipate that it would not
satisfy the Open Source Definition (OSD) as currently
interpreted~\cite{perens1999osi}. The OSD requires that licenses not
discriminate against ``fields of endeavour'' (OSD criterion 6), and the
consent layer's distinction between interpretive and non-interpretive
use could be characterised as discriminating against AI training.

We regard this as a limitation of the OSD, not of the PMPL. The OSD was
drafted in 1998, when the idea of non-interpretive use did not exist.
Its ``fields of endeavour'' criterion was designed to prevent licenses
from restricting use to, say, non-military purposes---not to prevent
licenses from distinguishing between human engagement and automated
extraction. We suggest that the OSD should be updated to recognise the
interpretive/non-interpretive distinction, but we do not condition the
PMPL's validity on this recognition.

% ============================================================================
% 9. IMPLEMENTATION
% ============================================================================
\section{Implementation}
\label{sec:implementation}

\subsection{SPDX Integration}

The PMPL-1.0 is identified by the SPDX License Identifier
\texttt{PMPL-1.0-or-later}. Source files include the standard SPDX
header:

\begin{lstlisting}
// SPDX-License-Identifier: PMPL-1.0-or-later
// SPDX-FileCopyrightText: 2026 Author Name
\end{lstlisting}

The ``-or-later'' suffix follows the convention established by the GPL
and LGPL, indicating that the work may be used under the specified
version or any later version published by the Palimpsest Stewardship
Council.

\subsection{Machine-Readable Consent Metadata}

The consent metadata specification defines a structured format for
expressing consent preferences. The canonical location is
\texttt{.palimpsest/consent.json} in the repository root, but the
specification also supports:

\begin{itemize}
    \item SPDX annotations in individual files
    \item HTTP headers for web-served content
    \item DNS TXT records for domain-level preferences
    \item \texttt{robots.txt}-style directives for web crawlers
\end{itemize}

The consent metadata schema supports the following fields:

\begin{lstlisting}
{
  "$schema": "https://palimpsest.dev/
    consent/v1.schema.json",
  "version": "1.0",
  "scope": "repository",
  "consent": {
    "interpretive": "granted",
    "non-interpretive": {
      "ai-training": "<policy>",
      "content-aggregation": "<policy>",
      "search-indexing": "<policy>",
      "academic-research": "<policy>",
      "automated-summarization": "<policy>"
    }
  },
  "conditions": {
    "attribution": true,
    "notification": false,
    "documentation": true,
    "lineage-preservation": true
  },
  "contact": "<email>",
  "updated": "<ISO-8601-date>",
  "signature": "<quantum-safe-sig>"
}
\end{lstlisting}

Where \texttt{<policy>} is one of: \texttt{"granted"} (permitted without
conditions), \texttt{"requires-notice"} (permitted with documentation
obligation), \texttt{"requires-permission"} (requires explicit
approval), or \texttt{"denied"} (not permitted under this license).

\subsection{Integration with AI Training Pipelines}

For the consent metadata to be effective, it must be discoverable by
AI training pipelines. We propose a multi-layered discovery mechanism:

\begin{enumerate}
    \item \textbf{Repository level}: The \texttt{.palimpsest/consent.json}
          file is checked during data collection.
    \item \textbf{Package level}: Package registries (npm, crates.io,
          PyPI) include consent metadata in package manifests.
    \item \textbf{Web level}: A \texttt{.well-known/palimpsest.json}
          endpoint serves consent metadata for web-hosted content.
    \item \textbf{DNS level}: A \texttt{\_palimpsest} TXT record
          provides domain-wide consent preferences.
\end{enumerate}

This layered approach ensures that consent metadata is discoverable
regardless of how the content is accessed. A training pipeline that
clones repositories checks the \texttt{.palimpsest/} directory. A
pipeline that scrapes websites checks the \texttt{.well-known/} endpoint.
A pipeline that queries package registries checks the manifest metadata.

\subsection{Provenance and Quantum-Safe Signatures}

The PMPL includes optional support for quantum-safe cryptographic
signatures on provenance metadata. This is forward-looking: current
cryptographic signatures (RSA, ECDSA) will become vulnerable when
large-scale quantum computers become available~\cite{nist-pqc}. The
PMPL supports three NIST-approved post-quantum algorithms: ML-DSA
(FIPS 204), SLH-DSA (FIPS 205), and FN-DSA (FIPS 206).

Provenance signatures serve two functions in the consent framework.
First, they authenticate the consent metadata itself, preventing
tampering (e.g., an AI company modifying a repository's consent file
to grant permission that was not given). Second, they create an
immutable audit trail, allowing retrospective verification of what
consent was in effect at any point in time.

\subsection{Tooling}

The reference implementation includes three Rust CLI tools:

\begin{itemize}
    \item \texttt{pmpl-sign}: Signs files and provenance metadata with
          quantum-safe signatures.
    \item \texttt{pmpl-verify}: Verifies signatures and provenance
          chains.
    \item \texttt{pmpl-audit}: Audits a repository for PMPL compliance,
          including consent metadata validation.
\end{itemize}

These tools are themselves licensed under the PMPL-1.0, ensuring that
the compliance infrastructure is subject to the same protections as the
works it protects.

% ============================================================================
% 10. DISCUSSION
% ============================================================================
\section{Discussion}
\label{sec:discussion}

\subsection{Implications for AI Governance}

The consent layer model has implications beyond software licensing. It
suggests a general principle: \emph{automated systems that process
creative works should be subject to different rules than humans who
engage with them}. This principle could inform:

\begin{itemize}
    \item \textbf{Regulatory frameworks}: The EU AI Act already
          distinguishes between different risk levels of AI systems.
          The interpretive/non-interpretive distinction could inform
          a more granular risk taxonomy.
    \item \textbf{Platform governance}: Social media platforms could
          distinguish between human sharing (interpretive) and automated
          scraping (non-interpretive) in their terms of service.
    \item \textbf{Academic publishing}: The ongoing debate about AI use
          in academic research could benefit from a clear distinction
          between using AI as a tool (interpretive) and using scholarly
          works as training data (non-interpretive).
\end{itemize}

\subsection{Open-Source Sustainability}

The open-source sustainability crisis is well-documented. Maintainers
burn out~\cite{eghbal2020working}. Critical infrastructure depends on
unpaid volunteers. Companies extract enormous value from open-source
software while contributing little back.

AI training exacerbates this dynamic. When a company trains a model on
open-source code and sells access to the model, the value chain runs
from open-source contributors (who created the training data) to the
model provider (who captured the value) to the end user (who pays for
access). The contributors receive nothing---not even the indirect benefit
of more developers learning from their code, since the model replaces
the learning process.

The consent layer does not solve the sustainability problem. But it
creates a \emph{leverage point}. If non-interpretive use requires notice
or consent, creators have a basis for negotiating compensation, imposing
conditions, or simply saying ``no.'' This is not anti-AI; it is
\emph{pro-autonomy}. The same principle underlies every labour
negotiation: the right to withhold one's labour (or, here, one's
creative output) is the foundation of fair exchange.

\subsection{Cultural Preservation in the Age of Generative AI}

Generative AI systems produce vast quantities of content that
statistically resembles human creative output. This creates a paradox:
the more AI-generated content exists, the harder it becomes to identify
and preserve authentic human creative expression. The emotional lineage
framework provides a mechanism for distinguishing between works with
genuine cultural roots and works that are statistically derived from
such roots.

The \texttt{LINEAGE.md} file, provenance signatures, and consent metadata
together create a \emph{cultural fingerprint}---a set of markers that
indicate not just who created a work, but \emph{why}, in \emph{what
context}, and with \emph{what meaning}. AI-generated content, by
definition, cannot produce authentic lineage documentation. The
distinction is not about technical sophistication; it is about
\emph{authenticity}---the relationship between a work and the human
experience that produced it.

\subsection{Objections and Responses}

\subsubsection{``This Is Anti-AI''}

The PMPL does not prohibit AI training. It requires transparency
(documentation of non-interpretive use) and provides a mechanism for
creators to express their preferences. This is no more ``anti-AI'' than
GDPR is ``anti-internet.'' Requiring consent for data processing is a
fundamental principle of information governance, not an act of
technological hostility.

\subsubsection{``Emotional Lineage Is Too Vague''}

The ``reasonable efforts'' standard, combined with the ``where
documented'' qualifier, grounds the obligation in verifiable facts. A
creator who documents their work's cultural context in
\texttt{LINEAGE.md} has created a concrete obligation. A creator who
does not has created no additional obligation beyond standard attribution.
The vagueness objection conflates the concept (broad) with the
obligation (specific).

\subsubsection{``This Will Fragment the Ecosystem''}

Any new license creates some fragmentation. The PMPL mitigates this by
building on the MPL-2.0, which is already widely understood and legally
tested. The additional provisions (emotional lineage, consent layer) are
\emph{additive}---they do not modify any MPL-2.0 permission or
obligation. Code licensed under the PMPL can be combined with MPL-2.0
code in the same project, minimising compatibility concerns.

\subsubsection{``Consent Metadata Will Be Ignored''}

It may be. Many Creative Commons metadata tags are ignored today. Many
\texttt{robots.txt} directives are violated. The existence of
non-compliance does not invalidate the legal instrument. Copyright itself
is routinely infringed; this does not make copyright law meaningless.
The consent layer provides a \emph{legal basis for action} when
non-compliance occurs, not a technological guarantee of compliance.

\subsection{Limitations and Future Work}

This paper and the PMPL-1.0 have several limitations that future work
should address:

\begin{itemize}
    \item \textbf{Enforcement infrastructure}: While the legal basis is
          sound, practical enforcement requires monitoring tools that
          can detect non-interpretive use of specific works. This is a
          technical challenge that current provenance tracking technology
          is only beginning to address.
    \item \textbf{International harmonisation}: Copyright and contract
          law vary by jurisdiction. While the PMPL's contractual approach
          provides broad enforceability, specific provisions may require
          adaptation for different legal systems.
    \item \textbf{Community adoption}: The success of any license depends
          on adoption. The PMPL is new and untested in court. Building
          the community and jurisprudential track record will take time.
    \item \textbf{The boundary problem}: The interpretive/non-interpretive
          distinction, while useful, has edge cases that will require
          ongoing refinement through community practice and legal
          interpretation.
    \item \textbf{Collective consent}: The current framework addresses
          individual creator consent. Works with many contributors raise
          collective action problems that the current governance
          structure (Stewardship Council) only partially addresses.
\end{itemize}

% ============================================================================
% 11. CONCLUSION
% ============================================================================
\section{Conclusion}
\label{sec:conclusion}

The open-source licensing frameworks that have served us for four decades
were built on an assumption that is no longer universal: that the primary
consumers of creative works are human beings capable of engaging with
their meaning. The emergence of non-interpretive AI systems---systems
that process works as statistical inputs without semantic
comprehension---creates a structural gap in every existing license.

This paper has introduced two concepts to address this gap: emotional
lineage protection, which formalises the obligation to preserve the
narrative, cultural, and symbolic meaning embedded in creative works; and
the consent layer for non-interpretive systems, which provides a
machine-readable mechanism for creators to express their preferences
about automated processing of their work.

These concepts are implemented in the Palimpsest-MPL License (PMPL-1.0),
which extends the Mozilla Public License 2.0 with ethical use
requirements, quantum-safe provenance, and the consent metadata
framework described herein. The PMPL is not anti-technology. It is
pro-consent. It applies the same principle that underlies privacy law,
labour law, and medical ethics: that the people affected by a system's
operation should have a voice in how that system uses their contributions.

The palimpsest metaphor reminds us that creative works are layered.
Beneath the visible surface---the current version, the compiled binary,
the rendered output---lie layers of history, meaning, and human
connection. A licensing framework adequate to the age of AI must protect
all of these layers, not just the ones that machines can see.

We invite the open-source community, legal scholars, and AI governance
researchers to engage with these ideas---to critique, refine, and build
upon them. Like the palimpsest itself, this work is a layer in an
ongoing conversation. It does not claim to be the final word. It claims
only to be a necessary one.

% ============================================================================
% ACKNOWLEDGMENTS
% ============================================================================
\section*{Acknowledgments}

The author thanks the Palimpsest Stewardship Council for their
contributions to the development of the PMPL-1.0 license text, and the
broader open-source licensing community for decades of foundational work
upon which this paper builds.

% ============================================================================
% REFERENCES
% ============================================================================
\begin{thebibliography}{99}

\bibitem{stallman1985gnu}
R.~Stallman, ``The GNU Manifesto,'' \emph{Dr.\ Dobb's Journal},
vol.~10, no.~3, pp.~30--35, 1985.

\bibitem{stallman2002free}
R.~M.~Stallman, \emph{Free Software, Free Society: Selected Essays of
Richard M.\ Stallman}. Boston, MA: GNU Press, 2002.

\bibitem{perens1999osi}
B.~Perens, ``The Open Source Definition,'' in \emph{Open Sources: Voices
from the Open Source Revolution}, C.~DiBona, S.~Ockman, and M.~Stone,
Eds. Sebastopol, CA: O'Reilly Media, 1999, pp.~171--188.

\bibitem{mongodb2018sspl}
MongoDB, Inc., ``Server Side Public License (SSPL) v1,'' 2018.
[Online]. Available: \url{https://www.mongodb.com/licensing/
server-side-public-license}

\bibitem{ehmke2019hippocratic}
C.~A.~Ehmke, ``The Hippocratic License: An Ethical License for Open
Source Projects,'' 2019.
[Online]. Available: \url{https://firstdonoharm.dev/}

\bibitem{anti996}
996.ICU Contributors, ``Anti 996 License Version 1.0,'' 2019.
[Online]. Available: \url{https://github.com/996icu/996.ICU}

\bibitem{copilot2022controversy}
A.~Butterick, ``GitHub Copilot Investigation,'' 2022.
[Online]. Available: \url{https://githubcopilotinvestigation.com/}

\bibitem{andersen2023}
S.~Andersen et~al.\ v.\ Stability AI Ltd.\ et~al., No.\ 3:23-cv-00201
(N.D.\ Cal.\ 2023).

\bibitem{nyt2023openai}
The New York Times Company v.\ Microsoft Corporation et~al.,
No.\ 1:23-cv-11195 (S.D.N.Y.\ 2023).

\bibitem{gplv3}
Free Software Foundation, ``GNU General Public License, Version 3,''
June 2007.
[Online]. Available: \url{https://www.gnu.org/licenses/gpl-3.0.html}

\bibitem{mit-license}
Open Source Initiative, ``The MIT License,'' n.d.
[Online]. Available: \url{https://opensource.org/licenses/MIT}

\bibitem{apache2}
Apache Software Foundation, ``Apache License, Version 2.0,'' January
2004.
[Online]. Available: \url{https://www.apache.org/licenses/LICENSE-2.0}

\bibitem{mpl2}
Mozilla Foundation, ``Mozilla Public License, Version 2.0,'' January
2012.
[Online]. Available: \url{https://www.mozilla.org/en-US/MPL/2.0/}

\bibitem{eu-ai-act-2024}
European Parliament and Council, ``Regulation (EU) 2024/1689 laying down
harmonised rules on artificial intelligence (AI Act),'' \emph{Official
Journal of the European Union}, L series, 2024.

\bibitem{eu-dsm-directive}
European Parliament and Council, ``Directive (EU) 2019/790 on copyright
and related rights in the Digital Single Market,'' \emph{Official Journal
of the European Union}, L 130, pp.~92--125, 2019.

\bibitem{samuelson1984copyright}
P.~Samuelson, ``CONTU Revisited: The Case Against Copyright Protection
for Computer Programs in Machine-Readable Form,'' \emph{Duke Law
Journal}, vol.~1984, no.~4, pp.~663--769, 1984.

\bibitem{clifford1988predicament}
J.~Clifford, \emph{The Predicament of Culture: Twentieth-Century
Ethnography, Literature, and Art}. Cambridge, MA: Harvard University
Press, 1988.

\bibitem{gdpr2016}
European Parliament and Council, ``Regulation (EU) 2016/679 on the
protection of natural persons with regard to the processing of personal
data (General Data Protection Regulation),'' \emph{Official Journal of
the European Union}, L 119, pp.~1--88, 2016.

\bibitem{lemley2023generative}
M.~A.~Lemley, ``How Broad is the Right to Train AI?,'' Working Paper,
Stanford Law School, 2023.

\bibitem{lessig2006code}
L.~Lessig, \emph{Code: And Other Laws of Cyberspace, Version 2.0}. New
York: Basic Books, 2006.

\bibitem{jacobsen2008}
Jacobsen v.\ Katzer, 535 F.3d 1373 (Fed.\ Cir.\ 2008).

\bibitem{farnsworth2004contracts}
E.~A.~Farnsworth, \emph{Farnsworth on Contracts}, 3rd ed. New York:
Aspen Publishers, 2004.

\bibitem{berne1886}
World Intellectual Property Organization, ``Berne Convention for the
Protection of Literary and Artistic Works,'' 1886, as amended 1979.

\bibitem{vara1990}
Visual Artists Rights Act of 1990, 17 U.S.C.\ \S\S~106A, 113(d)
(1990).

\bibitem{unesco2003}
UNESCO, ``Convention for the Safeguarding of the Intangible Cultural
Heritage,'' 2003.

\bibitem{wipo-tk}
World Intellectual Property Organization, ``Intergovernmental Committee
on Intellectual Property and Genetic Resources, Traditional Knowledge
and Folklore,'' established 2000.

\bibitem{nist-pqc}
National Institute of Standards and Technology, ``Post-Quantum
Cryptography Standardization,'' FIPS 204, 205, 206, 2024.

\bibitem{cc-by}
Creative Commons, ``Attribution 4.0 International (CC BY 4.0),'' 2013.
[Online]. Available: \url{https://creativecommons.org/licenses/by/4.0/}

\bibitem{rail2022}
BigScience, ``The RAIL License Family,'' 2022.
[Online]. Available: \url{https://www.licenses.ai/}

\bibitem{eghbal2020working}
N.~Eghbal, \emph{Working in Public: The Making and Maintenance of Open
Source Software}. San Francisco, CA: Stripe Press, 2020.

\bibitem{us-copyright-act}
Copyright Act of 1976, 17 U.S.C.\ \S\S~101--810.

\bibitem{davidson2004click}
A.~Davidson, ``Shrinkwrap and Clickwrap Agreements,'' \emph{St.\ John's
Law Review}, vol.~78, no.~3, pp.~613--638, 2004.

\bibitem{ftc-act}
Federal Trade Commission Act, 15 U.S.C.\ \S\S~41--58.

\bibitem{netz2007archimedes}
R.~Netz and W.~Noel, \emph{The Archimedes Codex: How a Medieval Prayer
Book Is Revealing the True Genius of Antiquity's Greatest Scientist}.
Philadelphia, PA: Da Capo Press, 2007.

\bibitem{lessig2004free}
L.~Lessig, \emph{Free Culture: How Big Media Uses Technology and the
Law to Lock Down Culture and Control Creativity}. New York: Penguin
Press, 2004.

\bibitem{benkler2006wealth}
Y.~Benkler, \emph{The Wealth of Networks: How Social Production
Transforms Markets and Freedom}. New Haven, CT: Yale University Press,
2006.

\bibitem{vaidhyanathan2001copyrights}
S.~Vaidhyanathan, \emph{Copyrights and Copywrongs: The Rise of
Intellectual Property and How It Threatens Creativity}. New York: New
York University Press, 2001.

\bibitem{zittrain2008future}
J.~Zittrain, \emph{The Future of the Internet---And How to Stop It}.
New Haven, CT: Yale University Press, 2008.

\bibitem{moglen2003dotcommunist}
E.~Moglen, ``The dotCommunist Manifesto,'' 2003.
[Online]. Available: \url{https://moglen.law.columbia.edu/
publications/dcm.html}

\bibitem{kelty2008two}
C.~M.~Kelty, \emph{Two Bits: The Cultural Significance of Free
Software}. Durham, NC: Duke University Press, 2008.

\bibitem{coyle2019value}
D.~Coyle, ``Do-it-yourself Digital: The Production Boundary, the
Productivity Puzzle and Economic Welfare,'' \emph{Economica}, vol.~86,
no.~344, pp.~750--774, 2019.

\bibitem{sobel2017artificial}
B.~L.~W.~Sobel, ``Artificial Intelligence's Fair Use Crisis,''
\emph{Columbia Journal of Law \& the Arts}, vol.~41, no.~1, pp.~45--97,
2017.

\bibitem{grimmelmann2016no}
J.~Grimmelmann, ``There's No Such Thing as a Computer-Authored Work---And
It's a Good Thing, Too,'' \emph{Columbia Journal of Law \& the Arts},
vol.~39, no.~3, pp.~403--416, 2016.

\bibitem{sag2019copyright}
M.~Sag, ``Copyright and Copy-Reliant Technology,'' \emph{Northwestern
University Law Review}, vol.~112, no.~6, pp.~1607--1682, 2019.

\end{thebibliography}

\end{document}
